\documentclass[conference]{IEEEtran}
\usepackage{cite}
\usepackage{amsmath,amssymb,amsfonts}
\usepackage{graphicx}
\usepackage{textcomp}
\usepackage{xcolor}
\usepackage{array}
\usepackage{url}
\usepackage{balance}

\newcommand{\designonly}{\textsc{Design-Only}}
\newcommand{\closure}{\textit{stereochemical closure}}
\newcolumntype{P}[1]{>{\raggedright\arraybackslash}p{#1}}

\begin{document}
\raggedbottom

\title{Why Does Life Require Chirality? An Evidence-Aware Proposal for Testing Stereochemical Closure Rather Than a Preferred Hand}

\author{\IEEEauthorblockN{Simulated Four-Agent Research Workflow}
\IEEEauthorblockA{Survey Agent $\rightarrow$ Idea Agent $\rightarrow$ ExperimentDesign Agent $\rightarrow$ Author Agent\\
Origin-of-Life Chemistry, Molecular Biology, and Scientific Evidence Modelling\\
Document state: \texttt{PROPOSAL\_NO\_OBSERVED\_RESULTS}}}

\maketitle

\begin{abstract}
The molecules of terrestrial life are chiral, but the popular phrase ``life requires one handedness'' hides several different claims. Proteins predominantly use L-amino acids, whereas the sugars in DNA and RNA have the D configuration; the relevant organization is biological homochirality and compatible stereochemical recognition, not a universal visual ``right hand.'' This paper asks whether chirality is functionally necessary for known chemical-life architectures and whether that statement can be generalized to all possible life. Evidence from template-directed RNA chemistry shows that matching stereochemical inputs can support extension whereas opposing inputs can inhibit it. Reviews and prebiotic-chemistry studies make several routes to enantioenrichment plausible, but do not identify a unique historical origin. A cross-chiral RNA polymerase ribozyme further shows that a common same-handed convention is not the only conceivable functional relation. We propose a \designonly{} four-agent research program centered on \closure: a reproducible compatibility relation among monomers, templates, catalysts, and recognition partners. A source-versioned evidence matrix compares homochiral, cross-chiral, mixed, and explicitly unsupported alternative architectures against preregistered capability criteria. The permitted conclusion is conditional support, alternative-architecture support, insufficient evidence, or definition-sensitive abstention. The project reports no synthesis, biological experiment, simulation result, mirror organism, or new observation.
\end{abstract}

\begin{IEEEkeywords}
chirality, homochirality, origin of life, RNA world, stereochemistry, molecular recognition, mirror chemistry, life definition, scientific evidence
\end{IEEEkeywords}

\section{Introduction}

Chirality is the property of an object or molecule whose mirror image cannot be superimposed on it. In a non-chiral physical environment, two enantiomers have the same bulk thermodynamic status but can behave differently when they interact with another chiral object. This becomes central in biology, where information-bearing polymers, folded catalysts, membranes, receptors, and substrates form geometrically constrained assemblies. A molecule that is the wrong mirror form may fail to fit, may change a conformation, or may be processed through a different reaction path. For this reason, chirality is often presented as a defining property of life.

That presentation needs qualification. It is inaccurate to say that DNA, RNA, and proteins all use a single arbitrary ``right-handed'' form. Terrestrial proteins predominantly use L-amino acids, whereas the sugars in DNA and RNA are D-configured. Their secondary and higher-order structures are also chiral, but the relevant observation is a coordinated stereochemical convention, not one label applied uniformly to all molecules. A scientific account must distinguish the chirality of individual components, homochirality within a polymer family, and compatibility relations across a biochemical network.

The stronger proposition---that all possible life must be chiral or homochiral---does not follow immediately from terrestrial biology. It depends on what counts as life, what chemical architectures are included in the reference class, and what level of stereochemical organization the proposition demands. NASA's frequently used working definition refers to ``a self-sustaining chemical system capable of Darwinian evolution'' \cite{nasa}; it does not specify a terrestrial L/D convention. A definition based on metabolism, autopoiesis, information processing, or adaptive evolution can impose a different evidentiary burden on a claim of necessity.

This report uses a simulated four-agent workflow to convert the broad question into a falsifiable research proposal. The Survey Agent constructs an evidence boundary. The Idea Agent searches alternatives and selects a closure-based hypothesis. The ExperimentDesign Agent specifies a source-bound, non-operational comparison. The Author Agent writes an IEEE-style proposal under fixed claim constraints. The central conclusion is deliberately conditional: homochirality is strongly functionally important in known biopolymer systems, but neither a particular terrestrial hand nor universal chirality of all conceivable life is established by the evidence reviewed here.

\section{What Is Actually Known About Biological Handedness}

\subsection{Molecular chirality and network homochirality}

At the molecular level, chiral compounds have enantiomers. At the biological level, a population of molecules can be enriched in one enantiomer, and a polymer can be assembled using a repeated stereochemical convention. The latter property is homochirality. It makes the geometry of a chain and the shape of its interaction surfaces more repeatable. The essential point is relational: a catalyst must recognize a substrate and a template must admit a compatible building block. The biological effect cannot be read from the isolated molecule alone.

Blackmond describes biological homochirality as a signature of life and surveys mechanisms by which an enantiomer could come to dominate from a presumably racemic prebiotic setting \cite{blackmond}. Hein and Blackmond likewise connect amino-acid and sugar enantioenrichment with the assembly of informational polymers \cite{hein}. These reviews support the existence of a major origin-of-life problem; they do not prove that one mechanism was historical or that all chemical life must follow terrestrial chemistry. That distinction is retained throughout this proposal.

\begin{table*}[!t]
\caption{Four claims commonly compressed into ``life requires chirality.''}
\label{tab:claims}
\centering
\renewcommand{\arraystretch}{1.20}
\begin{tabular}{|P{0.19\textwidth}|P{0.26\textwidth}|P{0.27\textwidth}|P{0.17\textwidth}|}
\hline
\textbf{Claim layer} & \textbf{Question} & \textbf{Evidence needed} & \textbf{It cannot establish alone} \\
\hline
Molecular chirality & Does a molecule have non-superimposable mirror forms? & Chemical structure and stereochemical designation. & That the molecule is biologically selected or necessary for life. \\
\hline
Biological homochirality & Does a polymer family or network use a consistent convention? & Composition and source-bound structural/functional evidence. & Why the convention originated or that its sign is universal. \\
\hline
Functional closure & Does a stated architecture preserve copying, folding, and recognition through compatible relations? & Direct function evidence plus defined participants and failure modes. & That every possible life architecture needs the same relation. \\
\hline
Universal necessity & Must every system counted as life be chiral or homochiral? & Explicit life definition, broad alternative-architecture comparison, and exclusion logic. & A conclusion from terrestrial examples only. \\
\hline
\end{tabular}
\end{table*}

\subsection{Terrestrial conventions are coordinated, not a single hand}

The conventional descriptors L and D are configurational conventions, not a general assertion that one enantiomer is visually left- or right-handed in every context. Proteins are overwhelmingly composed of L-amino acids; the nucleic acids use D-ribose or D-deoxyribose residues. This mixed wording is important because it prevents an error at the first step of the question. A biochemical network is not uniform because every constituent has one global sign. It is uniform because the stereochemical relations among its constituents are compatible and reproducible.

That compatibility affects both information and function. An informational polymer must form a repeatable interaction with a template or partner, while a functional polymer must fold into a repeatable structure and present interaction surfaces. The same physical property can therefore appear as a replication constraint, a structural constraint, and a recognition constraint. A future origin-of-life system may organize these functions differently; a report must state which of these capability requirements it assumes rather than using ``life'' as an undefined shortcut.

\subsection{Why template chemistry matters}

Joyce \textit{et al.} reported a landmark experiment in which poly(C)-directed oligomerization of activated guanosine mononucleotides proceeded readily when monomers matched the template's handedness and was much less efficient for the opposite handedness \cite{joyce1984}. In racemic mixtures, opposite-handed monomers could be incorporated as chain terminators. The result is direct evidence for a chirality-dependent bottleneck in a specific template-directed RNA-like setting. It does not show that every historical prebiotic polymer behaved this way, nor that any future chemical life must use exactly this chemistry.

The conditional scope matters. A matched convention can be selected because it improves the continuity of a template process. If continuity makes longer informational polymers possible, it can feed back into function and selection. Chen and Ma modeled a route by which polymer-level chiral selection, template copying, and amplification could generate a strong deviation from a racemic monomer pool \cite{chenma}. Their work motivates an evolutionary hypothesis, but it remains a model with assumptions about the RNA-world setting. Its output is neither a recovered early-Earth observation nor a proof of an inevitable origin.

\begin{figure*}[!t]
\centering
\fbox{\parbox{0.94\textwidth}{\centering
\vspace{4pt}
\textbf{Evidence-bound four-agent flow}\\[5pt]
\begin{tabular}{c@{\hspace{8pt}$\longrightarrow$\hspace{8pt}}c@{\hspace{8pt}$\longrightarrow$\hspace{8pt}}c@{\hspace{8pt}$\longrightarrow$\hspace{8pt}}c}
\fbox{\parbox{0.15\textwidth}{\centering \textbf{Survey}\\Chirality layers\\Evidence plan\\Gap ledger}} &
\fbox{\parbox{0.15\textwidth}{\centering \textbf{Idea}\\Alternatives\\Closure hypothesis\\Reject preferred hand}} &
\fbox{\parbox{0.18\textwidth}{\centering \textbf{ExperimentDesign}\\Evidence matrix\\Architecture states\\Abstention}} &
\fbox{\parbox{0.15\textwidth}{\centering \textbf{Author}\\IEEE proposal\\Claim audit\\No results claim}}
\end{tabular}\\[7pt]
\fbox{\parbox{0.80\textwidth}{\centering \textbf{Cross-stage invariant:} no model or review output proves a universal law of life, reports a new biological result, or demonstrates a mirror organism.}}
\par\vspace{4pt}}}
\caption{Claim and evidence flow for the simulated research process. An editable Draw.io version is included in the accompanying artifact directory.}
\label{fig:agents}
\end{figure*}

\section{Survey Synthesis: Mechanisms, Boundaries, and Gaps}

\subsection{Plausible routes to enantioenrichment}

The origin of homochirality is not settled by observing homochirality today. Chemical and physical mechanisms can introduce a small bias, select a crystal form, amplify an initial imbalance, or couple one class of biomonomers to another. The literature contains many such proposals, including autocatalytic amplification, phase behavior, crystallization, light-mediated mechanisms, and surface effects. Sallembien \textit{et al.} organize chemical and physical scenarios at the levels of building blocks, functional polymers, and more elaborate chemical systems \cite{sallembien}. The multiplicity is a reason for comparative analysis rather than a reason to declare the problem solved.

Ozturk \textit{et al.} provide a useful primary example. They reported enantioenrichment of racemic ribo-aminooxazoline on magnetite followed by homochiral crystallization, and interpret the result as a plausible pathway to system-level homochirality under a proposed early-Earth setting \cite{ozturk}. The report supports the limited statement that a racemic precursor can be driven toward a homochiral product through a demonstrated physical-chemical process. It does not identify the actual origin event, establish that magnetite was necessary, or connect that product automatically to self-sustaining evolution.

\subsection{From polymer properties to life capabilities}

There is a scale problem between an enantiomeric excess and life. A selective crystallization result is not a functioning genetic system. A polymerization result is not open-ended Darwinian evolution. A folded protein library is not a cell. Skolnick, Zhou, and Gao used computational protein libraries to compare homochiral and mixed-chirality structural and functional properties, finding lower stability for their demi-chiral protein set while also reporting some residual fold and active-site-like features \cite{skolnick}. This is informative precisely because it does not endorse a binary story: a mixed set may retain some local function while lacking the stability or network properties required by a more demanding capability definition.

The Survey Agent therefore treats life capabilities as a ladder rather than a label. It asks whether a source demonstrates template continuity, reproducible structure, selective recognition, heritable information, or support for open-ended evolution. A record must be allowed to support a lower-level capability without being promoted to a higher-level conclusion. This avoids a common transition error in origin-of-life argumentation.

\begin{table}[!t]
\caption{Source roles and claim limits used by the Survey Agent.}
\label{tab:source-roles}
\centering
\renewcommand{\arraystretch}{1.16}
\begin{tabular}{|P{0.23\columnwidth}|P{0.25\columnwidth}|P{0.34\columnwidth}|}
\hline
\textbf{Source class} & \textbf{Registered role} & \textbf{Claim limit} \\
\hline
Primary template experiment & Demonstrate a stated stereochemical effect in a stated reaction context. & No universal life inference. \\
\hline
Primary enrichment experiment & Demonstrate a route to enantioenrichment or crystallization. & No historical-origin or evolvability inference. \\
\hline
Cross-chiral experiment & Demonstrate a defined alternative compatibility relation. & No mirror organism or general-life claim. \\
\hline
Computational model & Explore conditional mechanism behavior. & No direct observation or parameter-free prediction. \\
\hline
Review or definition source & Provide synthesis, vocabulary, or scope. & Cannot populate direct experimental capability endpoints. \\
\hline
\end{tabular}
\end{table}

\subsection{A counterexample that changes the hypothesis}

Sczepanski and Joyce reported a cross-chiral RNA polymerase ribozyme in which a D-RNA enzyme catalyzed templated polymerization of L-RNA components, with analogous enantiomeric behavior reported in the reverse relation \cite{sczepanski}. This result is not a denial of chirality. The system is intensely stereochemical. It is instead a counterexample to a less careful premise: that an informational system can function only if catalyst, template, and substrates all share a single same-handed convention.

The right response is to replace ``one shared hand'' with a more general compatibility question. What relation among participants preserves the relevant information and function? A homochiral relation is one answer. A cross-chiral relation can be another answer. An explicitly achiral architecture may be hypothesized, but cannot be called a counterexample until its required capabilities are independently supported. The Survey Agent stores this as an alternative-architecture gap rather than as a victory for any speculative mirror-life account.

\subsection{Research gaps accepted for downstream search}

The final Survey handoff accepts four gaps. First, a function-versus-universality gap: terrestrial observations are routinely stated as a law of all life without an explicit reference class. Second, a scale-bridging gap: enantioenrichment, polymer function, heredity, and evolution are not consistently linked by source-traceable evidence. Third, an alternative-architecture gap: cross-chiral interactions are known in limited systems but poorly incorporated into arguments about necessity. Fourth, a definition-sensitive inference gap: claims about life often omit the capabilities their definition requires. These gaps constrain the subsequent idea search.

\begin{table*}[!t]
\caption{Accepted research gaps and downstream commitments.}
\label{tab:gaps}
\centering
\renewcommand{\arraystretch}{1.18}
\begin{tabular}{|P{0.15\textwidth}|P{0.26\textwidth}|P{0.29\textwidth}|P{0.18\textwidth}|}
\hline
\textbf{Gap} & \textbf{Problem} & \textbf{Downstream commitment} & \textbf{Prohibited shortcut} \\
\hline
GAP-01 & Functional observations are promoted to a universal law. & State a life definition and architecture before testing necessity. & ``All life is chiral because Earth life is.'' \\
\hline
GAP-02 & Evidence is fragmented across chemical scales. & Record capability level and source role for every claim. & Treat enrichment as a demonstration of life. \\
\hline
GAP-03 & Alternatives to shared handedness are neglected. & Compare homochiral and cross-chiral closure explicitly. & Call a ribozyme a mirror organism. \\
\hline
GAP-04 & Life definitions are left implicit. & Run definition-sensitive analyses and preserve abstention. & Infer an achiral counterexample from a change of wording. \\
\hline
\end{tabular}
\end{table*}

\section{Idea Portfolio and Selected Hypothesis}

The Idea Agent compared four directions. D1, a historical atlas of proposed symmetry-breaking mechanisms, would be useful but could not discriminate functional necessity from origin narratives. D2, an operational test of \closure, separates compatible network relations from a preferred sign and directly addresses all four accepted gaps. D3, a mirror-chemistry feasibility argument, is scientifically valuable as a counterexample module but cannot become a claim that mirror life exists. D4, an extrapolation from terrestrial L-amino acids and D-sugars to every possible life, was rejected because it changes a contingent observation into a universal premise.

The selected D2 hypothesis is: \emph{for a stipulated chemical-life architecture that uses chiral informational or functional polymers, durable heredity and selective function require a reproducible stereochemical closure among participating monomers, templates, catalysts, and recognition partners; homochirality is the most strongly evidenced terrestrial realization, while the preferred sign and universal necessity remain unproven.} The wording deliberately includes a reference class. If the architecture changes, the claim must be re-evaluated rather than carried forward as an axiom.

The idea search also identifies failure modes. Closure could be an over-engineered name for a proxy that adds no explanatory value beyond existing functional descriptors. The literature may be too heterogeneous to compare. Cross-chiral records could show that the proposed categories are wrong rather than merely that same-handedness is too narrow. Or a plausible achiral architecture might satisfy the pre-registered capability criteria. All of these outcomes are informative and none would authorize a claim that a new form of life has been made.

\begin{figure}[!t]
\centering
\fbox{\parbox{0.90\columnwidth}{\centering
\textbf{Capability ladder and claim boundary}\\[4pt]
\fbox{Chiral molecule} $\rightarrow$ \fbox{Compatible polymer interaction} $\rightarrow$ \fbox{Heredity-supporting network} $\rightarrow$ \fbox{Life under declared definition}\\[5pt]
\parbox{0.82\columnwidth}{\centering Evidence may support one rung without supporting the next. A cross-chiral example can satisfy a compatibility relation without proving mirror life; an achiral proposal remains unclassified until it supports the specified capabilities.}}}
\caption{The proposal treats life-capability claims as a ladder. It prohibits promotion of lower-rung evidence to a universal life conclusion.}
\label{fig:ladder}
\end{figure}

\section{ExperimentDesign: A Conditional Closure Protocol}

\subsection{Research object and record contract}

The object of analysis is a time-versioned, source-bound evidence corpus, not a wet-lab system. A record may represent a primary experiment, a computational model, a review claim, a defined life criterion, a documented limitation, or an explicitly excluded speculative proposal. Every eligible record carries a stable identifier, source and publication date, exact supporting passage, source type, architecture state, capability endpoint, limitation, and claim-language state. The evidence corpus is designed so a reviewer can identify what a citation establishes and what it does not.

The architecture state is coded as \texttt{homochiral closure}, \texttt{cross-chiral closure}, \texttt{mixed unclosed}, or \texttt{achiral or unknown}. These are not biological labels. They are source-bound descriptions of the relation stated or demonstrated in a record. ``Unknown'' is a substantive value: a source that lacks an adequate stereochemical description cannot be silently assigned to a favorable group. A review may be used to discover a source or define a vocabulary, but it cannot by itself populate a direct experimental capability endpoint.

\subsection{Closure representation}

For a source record $r$, an architecture descriptor $A$, and a declared definition-of-life contract $D$, the proposal represents the closure assessment as
\begin{equation}
\mathcal{C}(r;A,D)=\big(t(r),s(r),q(r),h(r),e(r),u(r)\big).
\label{eq:closure}
\end{equation}
Here $t$ is evidence for template continuity, $s$ for reproducible structure, $q$ for selective recognition, $h$ for heritable information, $e$ for support of open-ended evolution, and $u$ records ambiguity, missingness, and source-model dependence. Equation \eqref{eq:closure} is not a probability that life exists and it does not measure intrinsic moral or biological value. It is a structured claim ledger.

The definition contract $D$ registers which capabilities are required before a system may count as life for the particular analysis. One analysis could use a Darwinian-evolution-oriented contract, another a metabolism/autopoiesis-oriented contract. A result is stable only if the conclusion survives the declared variation, or if the report precisely names the definition under which it holds. The design therefore treats definition sensitivity as a result, not a defect to hide.

\begin{table*}[!t]
\caption{Variables and claim controls for the closure protocol.}
\label{tab:variables}
\centering
\renewcommand{\arraystretch}{1.18}
\begin{tabular}{|P{0.18\textwidth}|P{0.17\textwidth}|P{0.31\textwidth}|P{0.15\textwidth}|}
\hline
\textbf{Field family} & \textbf{Role} & \textbf{Definition and control} & \textbf{Proposal-time state} \\
\hline
Record provenance & Evidence boundary & DOI/URL, date, source type, passage, independence relation, exclusion and ambiguity reason. & Required; no source invention. \\
\hline
Architecture state $A$ & Main comparison & Homochiral, cross-chiral, mixed-unclosed, or achiral/unknown as the source supports. & Planned coding only. \\
\hline
Capability vector & Functional context & $t,s,q,h,e$ are source-demonstrated, source-argued, not assessed, or out of scope. & Planned extraction only. \\
\hline
Definition contract $D$ & Scope control & Explicit life criterion, required capabilities, and permissible conclusion language. & Mandatory preregistration. \\
\hline
Uncertainty $u$ & Abstention control & Source ambiguity, model dependence, reviewer disagreement, and out-of-domain flag. & Mandatory. \\
\hline
Claim state & Governance control & Conditional support, alternative support, insufficient evidence, or definition-sensitive abstention. & Permitted output only. \\
\hline
\end{tabular}
\end{table*}

\subsection{Evidence states and prohibited outputs}

The only permitted top-level outputs are \texttt{closure functionally supported}, \texttt{alternative architecture supported}, \texttt{evidence insufficient}, and \texttt{definition-sensitive abstain}. They describe the corpus under its registered definition and source boundary. The prohibited labels are \texttt{life created}, \texttt{universal law proved}, and \texttt{mirror organism exists}. This language rule is as important as the calculation because a categorized evidence record can otherwise be mistaken for an experimental verdict about life.

\section{Validation Plan and Decision Rules}

\subsection{Corpus construction and review}

Before any synthesis, the protocol freezes inclusion rules, source window, architecture codebook, endpoint definitions, life-definition contracts, reviewer qualifications, adjudication process, and sensitivity analyses. Two domain-qualified reviewers independently annotate a pilot subset. Each non-missing label requires a passage anchor and a rationale. Disagreements can be adjudicated, retained as competing labels, or marked ineligible. A high agreement number is not the sole target; a field that consistently resists reliable coding may reveal that the proposed concept is too coarse or that public evidence cannot support the requested distinction.

Source dependence is recorded. Multiple papers may repeat a common experiment or source a claim from a single review. Such papers are not counted as independent confirmations merely because their titles differ. Later reviews may summarize earlier evidence, but their hindsight cannot be inserted into an earlier record. The corpus retains negative, ambiguous, and excluded records to prevent a success-only reconstruction of the origin-of-life literature.

\subsection{Comparators and endpoints}

The primary comparator contrasts records supporting a registered closure relation with records explicitly characterized as mixed and unclosed. The comparison is descriptive and evidence-weighted, not a claim of a pooled causal effect. The secondary comparator asks whether a cross-chiral closure can satisfy the same capability fields by a different relation. A third comparator preserves explicitly achiral proposals as hypotheses or unknowns until direct evidence supports the registered endpoints.

The protocol measures endpoint coverage and evidentiary consistency. For an endpoint set $K_D$ required under definition contract $D$, a simple transparent coverage measure is
\begin{equation}
\operatorname{Cov}(r;D)=\frac{\sum_{k\in K_D} w_k I_k(r)}{\sum_{k\in K_D} w_k},
\label{eq:coverage}
\end{equation}
where $w_k$ are preregistered importance weights and $I_k(r)$ equals one only when record $r$ has source-demonstrated support for capability $k$. Equation \eqref{eq:coverage} is a reporting device, not a proof that a system is alive. It cannot replace qualitative reading of source limitations, and it is not evaluated in this proposal.

\begin{table}[!t]
\caption{Decision rules for a future evidence synthesis.}
\label{tab:decisions}
\centering
\renewcommand{\arraystretch}{1.15}
\begin{tabular}{|P{0.27\columnwidth}|P{0.26\columnwidth}|P{0.29\columnwidth}|}
\hline
\textbf{Future evidence pattern} & \textbf{Permitted conclusion} & \textbf{Required action} \\
\hline
Closure records consistently support registered capabilities under a stated $D$. & Functional support for closure in that architecture class. & Report source roles, gaps, alternatives, and definition scope. \\
\hline
Cross-chiral records meet equivalent capabilities. & A same-handed convention is not the only compatible relation. & Preserve the distinction from a mirror-organism claim. \\
\hline
Mixed or achiral alternative meets declared criteria with direct evidence. & Alternative architecture support. & Re-evaluate the closure hypothesis and publish the counterexample. \\
\hline
Reviewer agreement or source coverage is weak. & Evidence insufficient or definition-sensitive abstention. & Preserve ambiguity ledger and avoid universal wording. \\
\hline
\end{tabular}
\end{table}

\subsection{Negative outcomes are informative}

The project must report a negative result if no reliable relationship appears between closure and the registered endpoints, if a definition contract changes the conclusion materially, or if cross-chiral cases cannot be represented without distortion. The research value then lies in identifying where the current language or evidence is inadequate. A negative result does not mean chirality has no relevance to known life; it means this particular generalization has not earned support. Conversely, a positive result under a chiral-polymer contract does not become a proof about every possible life.

\section{Alternative Architectures and the Scope of Universality}

\subsection{Mirror-related systems}

Mirror chemistry is a conceptual stress test. If an entire compatible biochemical network were reflected, its components could in principle preserve internal relations while being incompatible with terrestrial partners. Such a possibility demonstrates why ``the L/D sign used by Earth life'' is a poor candidate for a universal necessity. However, a conceptual mirror relation is not evidence that a self-sustaining mirror organism exists, and it is not a proposal to create one. The present report does not specify organisms, molecular sequences, laboratory protocols, or interventions.

The cross-chiral ribozyme result is more precise than a general mirror-life slogan. It shows that a defined chiral enzyme can operate on the opposite enantiomeric class under a particular molecular architecture \cite{sczepanski}. The result changes the hypothesis from shared handedness to closure relation. It also sets a methodological rule: a counterexample must preserve its exact scope. A ribozyme does not automatically provide compartmentalization, metabolism, heredity across generations, or open-ended evolution.

\subsection{Could life be achiral?}

An achiral molecule or material can participate in a chiral environment, and a network of individually achiral components may produce chiral assemblies or dynamical states. Therefore ``achiral life'' can mean several things: achiral monomers, no global enantiomeric excess, no stereochemical selectivity, or an architecture whose functional units acquire chirality only at a higher scale. These are not interchangeable. Any universal claim must classify them separately.

No source in the present evidence plan is treated as demonstrating an achiral system that fulfills a demanding chemical-life contract. The correct status is not impossibility but unclassified or unsupported. The protocol gives such a proposal a route to challenge the hypothesis: it must provide source-bound support for the same required capabilities, with adequate evidence for persistence, information, selective function, and the declared definition. This prevents both an argument from ignorance and unearned speculation.

\subsection{Definitions determine the burden of proof}

If life is defined minimally as a self-sustaining chemical system capable of Darwinian evolution, a candidate analysis must explain how heredity and variation persist. If life is defined using autopoiesis or metabolism, the emphasis shifts to boundary maintenance and network organization. In every case, chirality becomes a claimed means to an end, not the end itself. A claim that it is necessary has to show why alternatives cannot meet the relevant capability bundle.

\begin{table*}[!t]
\caption{Architecture classes and interpretation boundaries.}
\label{tab:architectures}
\centering
\renewcommand{\arraystretch}{1.17}
\begin{tabular}{|P{0.19\textwidth}|P{0.22\textwidth}|P{0.28\textwidth}|P{0.20\textwidth}|}
\hline
\textbf{Architecture class} & \textbf{Relation being tested} & \textbf{What current evidence can say} & \textbf{What it cannot say} \\
\hline
Homochiral closure & Compatible common convention within key polymer interactions. & Strong functional relevance for several terrestrial-like template, fold, and recognition systems. & That the terrestrial sign or exact chemistry is universal. \\
\hline
Cross-chiral closure & Complementary relation across opposite conventions. & Defined RNA-ribozyme evidence shows a nontrivial alternative to simple shared handedness. & That a complete mirror life form exists. \\
\hline
Mixed unclosed & No source-supported compatibility relation across interacting components. & Can reveal inhibition, reduced continuity, or unstable structure in a stated context. & That every mixed system is nonliving. \\
\hline
Achiral or unknown & Explicitly achiral proposal or insufficient stereochemical evidence. & Identifies an evidence obligation and possible counterexample space. & That achiral life is established or impossible. \\
\hline
\end{tabular}
\end{table*}

\section{Validity Threats, Governance, and Safety Boundaries}

\subsection{Construct and evidence validity}

The central construct threat is treating an evidence code as the phenomenon itself. Closure is a proposed organizing concept, not a measurable substance. A high coverage score can be obtained from weakly related papers if the source boundary is loose; a low score can arise because a capability was not measured rather than because it failed. The design controls this by recording source passages, source role, missingness, and claim limits. It also separates source-demonstrated, source-argued, and unassessed fields.

Temporal and publication bias matter. Successful or striking mechanisms are more likely to be published, reviewed, and cited than unproductive paths. A recent review can make a historical story appear more coherent than the earlier evidence warranted. The protocol prevents this by preserving publication dates, evidence provenance, exclusions, and independent-origin relationships. It does not claim to eliminate missing negative results.

\subsection{Language and governance boundary}

The term ``life'' has scientific, philosophical, and societal significance. A model cannot settle its definition through a numeric score. Every conclusion must state its definition contract, source set, architecture class, and uncertainty. The Author Agent must reject sentences that convert a functional observation into a metaphysical necessity, or a conceptual mirror relation into a report of a mirror organism. The literature includes primary experiments, models, and reviews with different roles; citations may not be used outside the role registered by the Survey Agent.

\subsection{Physical-work exclusion}

This is not a molecular-biology or origin-of-life laboratory protocol. It gives no instructions for synthesizing chiral molecules, constructing RNA, conducting directed evolution, modifying organisms, culturing cells, or making mirror biological systems. Any future physical investigation must be independently designed by qualified researchers under local chemical, biological, institutional, and ethical requirements. The current contribution is restricted to evidence organization, hypothesis formation, and conditional research design.

\begin{table}[!t]
\caption{Risk register and mandatory response.}
\label{tab:risks}
\centering
\renewcommand{\arraystretch}{1.15}
\begin{tabular}{|P{0.22\columnwidth}|P{0.20\columnwidth}|P{0.28\columnwidth}|}
\hline
\textbf{Risk} & \textbf{Failure signal} & \textbf{Required response} \\
\hline
Sign confusion & Text calls all biological molecules ``right-handed.'' & Restore L-amino-acid/D-sugar distinction and network framing. \\
\hline
Scale jump & Enrichment evidence is stated as life creation. & Downgrade to its demonstrated chemical capability. \\
\hline
Alternative erasure & Cross-chiral evidence is omitted or relabeled homochiral. & Retain separate architecture state and sensitivity analysis. \\
\hline
Definition drift & Result changes after life criteria change informally. & Register a new definition contract and report sensitivity. \\
\hline
Operational drift & Work turns into a synthesis or organism-engineering plan. & Stop and require qualified, authorized external process. \\
\hline
\end{tabular}
\end{table}

\section{Expected Contributions and Interpretation of Outcomes}

The first expected contribution is a vocabulary that prevents a molecular property, a terrestrial regularity, a functional condition, and a universal law from being used interchangeably. The second is a source-bound evidence matrix that preserves scale, uncertainty, and alternative architectures. The third is a conditional closure hypothesis that can fail: an alternative architecture may satisfy the same capability contract, or the evidence may not support a stable mapping. The fourth is a reporting discipline that makes definition sensitivity visible.

A future positive result would be modest. It would mean that, under a declared chiral-polymer chemical-life contract and a frozen evidence corpus, compatible stereochemical closure has stronger support than mixed-unclosed alternatives for the registered capabilities. It would not mean that Earth biology's sign is preferred by a universal law, that all life must be homochiral, or that an origin mechanism has been identified. A future alternative-architecture result would be equally valuable: it would refine or limit the hypothesis.

The expected public-facing answer is therefore two-part. Why does known life need chirality? Because its molecular machinery needs reproducible three-dimensional compatibility for copying, folding, catalysis, and recognition. Must every possible life be chiral in the same way? The current evidence does not establish that. The more testable claim is that a life-like chemical architecture needs a stable information-preserving compatibility structure; homochirality is a prominent terrestrial realization, not yet a definition-independent law.

\balance
\section{Minimum Preregistration and Evidence-Record Contract}

Before any future implementation, the project must register the source window, search terms, inclusion/exclusion logic, life-definition contracts, architecture vocabulary, capability endpoints, reviewer protocol, independence criteria, missing-data treatment, sensitivity analyses, decision rules, and permitted conclusion language. A change after reading later evidence creates a new analysis rather than a silent revision. The purpose is not bureaucracy; it is to prevent a desired conclusion from selecting its own definition and source set.

At the record level, the archive must preserve a stable identifier, citation, source passage, source role, architecture state, capability endpoint, limitation, uncertainty flag, relation to other records, and exact report-language state. The record must retain negative and inconclusive entries. A cross-chiral paper must not be transformed into a homochiral paper for convenience, and a review statement must not be transformed into a direct observation. These rules make the evidence auditable by both molecular-science and origin-of-life specialists.

\begin{table}[!t]
\caption{Minimum evidence record for a stereochemical-closure assessment.}
\label{tab:record}
\centering
\renewcommand{\arraystretch}{1.14}
\begin{tabular}{|P{0.28\columnwidth}|P{0.53\columnwidth}|}
\hline
\textbf{Record block} & \textbf{Required fields} \\
\hline
Registration & Life-definition contract, source window, codebook version, endpoints, review roles, sensitivity plan, and decision rule. \\
\hline
Provenance & Stable record ID, DOI/URL, date, source type, passage anchor, rights/access note, and dependence relation. \\
\hline
Architecture & Homochiral, cross-chiral, mixed-unclosed, achiral/unknown, stated assumptions, and exclusion rationale. \\
\hline
Capability evidence & Template, structure, recognition, heredity, evolution-support state, source role, and limitation. \\
\hline
Claim audit & Allowed wording, uncertainty, definition sensitivity, counterexample status, and downgrade/negative-result record. \\
\hline
\end{tabular}
\end{table}

\section{Conclusion}

Life as we know it is chiral because its molecular machines operate through geometrically specific relationships. Homochirality helps make those relationships repeatable across informational polymers, folded catalysts, and recognition systems. The evidence from template-directed RNA chemistry, molecular-function studies, and prebiotic enantioenrichment research supports a strong functional role for compatible stereochemistry in terrestrial-like chemical architectures. It does not demonstrate a single historical origin route, a globally preferred terrestrial hand, a mirror organism, or an impossibility theorem for achiral life.

This paper offers a more precise research program. It replaces a vague universal statement with a test of \closure{} under declared capability and life-definition contracts. It requires source provenance, alternative-architecture comparison, abstention, and negative-result reporting. The resulting conclusion can be scientifically useful even when it is narrow: homochirality is a powerful and evidenced solution to the problem of chemical compatibility; whether all possible life must instantiate chirality or a related closure is a remaining empirical and conceptual question.

\section*{Acknowledgment}

This document is a simulated four-agent research proposal. It reports no molecular synthesis, biological experiment, origin-of-life observation, computational result generated by this workflow, mirror organism, or universal proof.

\begin{thebibliography}{00}

\bibitem{blackmond} D. G. Blackmond, ``The Origin of Biological Homochirality,'' \emph{Cold Spring Harbor Perspectives in Biology}, vol. 2, Art. no. a002147, 2010, doi: 10.1101/cshperspect.a002147.

\bibitem{joyce1984} G. F. Joyce, G. M. Visser, C. A. A. van Boeckel, J. H. van Boom, L. E. Orgel, and J. van Westrenen, ``Chiral selection in poly(C)-directed synthesis of oligo(G),'' \emph{Nature}, vol. 310, pp. 602--604, 1984, doi: 10.1038/310602a0.

\bibitem{chenma} Y. Chen and W. Ma, ``The origin of biological homochirality along with the origin of life,'' \emph{PLoS Computational Biology}, vol. 16, Art. no. e1007592, 2020, doi: 10.1371/journal.pcbi.1007592.

\bibitem{sczepanski} J. T. Sczepanski and G. F. Joyce, ``A cross-chiral RNA polymerase ribozyme,'' \emph{Nature}, vol. 515, pp. 440--442, 2014, doi: 10.1038/nature13900.

\bibitem{ozturk} S. F. Ozturk, Z. Liu, J. D. Sutherland, and D. Sasselov, ``Origin of biological homochirality by crystallization of an RNA precursor on a magnetic surface,'' \emph{Science Advances}, vol. 9, Art. no. eadg8274, 2023, doi: 10.1126/sciadv.adg8274.

\bibitem{hein} J. E. Hein and D. G. Blackmond, ``On the Origin of Single Chirality of Amino Acids and Sugars in Biogenesis,'' \emph{Accounts of Chemical Research}, vol. 45, pp. 2045--2054, 2012, doi: 10.1021/ar200316n.

\bibitem{skolnick} J. Skolnick, H. Zhou, and M. Gao, ``On the possible origin of protein homochirality, structure, and biochemical function,'' \emph{Proceedings of the National Academy of Sciences}, vol. 116, pp. 26571--26579, 2019, doi: 10.1073/pnas.1908241116.

\bibitem{sallembien} Q. Sallembien, L. Bouteiller, J. Crassous, and M. Raynal, ``Possible chemical and physical scenarios towards biological homochirality,'' \emph{Chemical Society Reviews}, vol. 51, pp. 3436--3476, 2022, doi: 10.1039/D1CS01179K.

\bibitem{nasa} NASA Astrobiology, ``About Life Detection,'' working definition of life. Available: \url{https://astrobiology.nasa.gov/research/life-detection/about/}. Accessed: Aug. 31, 2026.

\end{thebibliography}
\end{document}
